\chapter{Cognitive impairment}

\medskip
\noindent\textit{\textbf{Under review.} This chapter is an AI-assisted educational draft; it carries no source citations and is not a substitute for professional medical advice.}

\section*{Definition and Overview}
Cognitive impairment is a broad term describing a decline in mental abilities such as memory, attention, language, problem-solving, and judgement. It ranges from mild changes that are noticeable to the person or family but do not interfere much with daily life, to severe impairment that makes independent living difficult. Cognitive impairment can be temporary or permanent, and can result from a wide range of causes, including normal ageing to a degree, medical illness, injury, medication, or brain disorders such as dementia. Distinguishing mild, age-related change from a more serious condition is an important task for health professionals, because some causes of cognitive impairment are treatable or reversible.

\section*{Epidemiology}
Cognitive impairment becomes more common with age, although it can affect younger people after head injury, stroke, infection, or with certain medical conditions. It is estimated that a significant proportion of older adults experience measurable cognitive decline, and the number of people living with dementia worldwide is growing as populations age. Mild cognitive impairment is more common than dementia and is diagnosed when memory or thinking problems are greater than expected for age but do not yet meet the threshold for dementia. Men and women are both affected, though the risk of some dementias differs between the sexes, partly because women tend to live longer.

\section*{Aetiology and Pathophysiology}
The causes of cognitive impairment are many and can be grouped into categories:

\begin{itemize}
    \item \textbf{Neurodegenerative disease:} conditions such as Alzheimer disease, vascular dementia, Lewy body dementia, and frontotemporal dementia, in which brain cells progressively lose function
    \item \textbf{Vascular causes:} strokes, small vessel disease, and reduced blood flow to the brain
    \item \textbf{Injury:} traumatic brain injury, including repeated concussions
    \item \textbf{Infection:} meningitis, encephalitis, or late-stage infections such as neurosyphilis
    \item \textbf{Metabolic and nutritional causes:} vitamin B12 deficiency, thyroid disease, liver or kidney failure
    \item \textbf{Medication and substances:} sedatives, anticholinergic drugs, alcohol, and recreational drugs
    \item \textbf{Delirium:} acute confusion caused by illness, surgery, infection, or dehydration
    \item \textbf{Psychiatric conditions:} severe depression can present with poor concentration and memory (so-called pseudodementia)
\end{itemize}

In neurodegenerative disease, abnormal proteins accumulate in the brain and disrupt communication between neurons, which gradually leads to cell death in regions important for memory and thinking.

\section*{Clinical Features}
\begin{itemize}
    \item Forgetfulness, especially of recent events, names, or appointments
    \item Difficulty finding words or following conversations
    \item Trouble with planning, problem-solving, or handling money
    \item Getting lost in familiar places or losing track of time
    \item Difficulty concentrating or being easily distracted
    \item Changes in mood, personality, or social behaviour
    \item Reduced judgement, such as poor decisions about safety or finances
    \item Difficulty with everyday tasks such as cooking, driving, or taking medication
\end{itemize}

\section*{Differential Diagnosis}
\begin{itemize}
    \item Normal ageing, where mild slowing and occasional forgetfulness occur without significant functional impact
    \item Mild cognitive impairment vs.\ dementia, distinguished mainly by the degree of effect on daily functioning
    \item Depression, anxiety, or stress, which can mimic memory and concentration problems
    \item Delirium, which has a sudden onset, fluctuates, and is often caused by an acute illness
    \item Medication side effects or substance misuse
    \item Vitamin deficiencies, thyroid disease, or other medical conditions
    \item Specific causes of dementia such as Alzheimer disease, vascular dementia, or frontotemporal dementia
\end{itemize}

\section*{When to Seek Medical Care}
See a doctor if you or a family member notice persistent or worsening problems with memory, thinking, or daily functioning, especially if the change comes on suddenly, follows a head injury, or is accompanied by confusion, weakness, or difficulty speaking.

\textbf{Call 000 for:}
\begin{itemize}
    \item Sudden confusion, weakness, or difficulty speaking (possible stroke)
    \item A sudden, severe change in consciousness or behaviour
    \item Collapse, a fit, or loss of consciousness
    \item Any sudden, severe symptom that worries you
\end{itemize}

\section*{Diagnosis and Investigations}
\begin{itemize}
    \item \textbf{History:} account from the person and a relative or carer about the onset, course, and impact of symptoms
    \item \textbf{Cognitive testing:} short screening tests such as the Mini-Mental State Examination or Montreal Cognitive Assessment
    \item \textbf{Detailed neuropsychological assessment:} for complex cases, to map strengths and weaknesses across different mental functions
    \item \textbf{Blood tests:} full blood count, thyroid function, vitamin B12, folate, glucose, and liver or kidney function to exclude reversible causes
    \item \textbf{Brain imaging:} CT or MRI to look for strokes, shrinkage of brain tissue, or other structural problems
    \item \textbf{Assessment of function:} how well the person manages daily activities, using reports from family or carers
    \item \textbf{Referral:} to a memory clinic, geriatrician, or neurologist for further assessment
\end{itemize}

\section*{Management}
\begin{itemize}
    \item \textbf{Treat reversible causes:} correcting vitamin deficiencies, adjusting medicines, treating thyroid disease or depression
    \item \textbf{Medication:} some medicines can slow the progression of certain dementias, though they are not a cure
    \item \textbf{Cognitive rehabilitation:} memory strategies, routines, and aids such as calendars, alarms, and labelled drawers
    \item \textbf{Occupational therapy:} support to maintain independence and adapt the home environment
    \item \textbf{Support services:} carer support, respite care, community nursing, and advocacy services
    \item \textbf{Safety planning:} review of driving, financial management, and falls risk
    \item \textbf{Regular follow-up:} monitoring of symptoms, function, and safety over time
\end{itemize}

\section*{Complications}
\begin{itemize}
    \item Progression of underlying disease, leading to increasing dependence on others
    \item Falls, fractures, and other injuries from impaired balance or judgement
    \item Wandering, getting lost, or leaving appliances on
    \item Poor nutrition, dehydration, and neglect of personal care
    \item Difficulty managing medication, leading to medical errors
    \item Depression, anxiety, and frustration in the affected person
    \item Carer stress, burnout, and family strain
\end{itemize}

\section*{Prognosis and Outcomes}
The outlook depends on the cause. Cognitive impairment due to a reversible cause, such as a vitamin deficiency or medication, may improve substantially or fully once the cause is treated. Mild cognitive impairment may remain stable, improve, or progress to dementia, with a substantial proportion of people developing dementia within a few years. Dementia is progressive, but the rate varies considerably, and many people live well for years with support, adaption, and treatment. Early assessment and planning help people make the most of their current abilities and arrange future care.

\section*{Prevention and Self-Care}
\begin{itemize}
    \item Stay physically active, which supports blood flow and brain health
    \item Keep mentally engaged with reading, puzzles, hobbies, and social activities
    \item Eat a balanced diet and control blood pressure, diabetes, and cholesterol
    \item Limit alcohol and avoid smoking
    \item Protect the head by wearing seatbelts and helmets
    \item Protect hearing and vision, as sensory loss is linked to cognitive decline
    \item Have regular health checks, especially with ageing
    \item Seek help early if memory or thinking problems arise
\end{itemize}
